% 強制回應開頭自我稱謂機制研究——會議論文初稿（中文版）
% 編譯：xelatex main-zh-TW.tex（需要 xeCJK）
\documentclass[10pt,twocolumn]{article}
\usepackage[a4paper,margin=2cm]{geometry}
\usepackage{xeCJK}
\setCJKmainfont[AutoFakeBold=2.5]{Microsoft JhengHei}
\usepackage{booktabs}
\usepackage{graphicx}
\usepackage{amsmath}
\usepackage{url}
\usepackage{hyperref}

\title{強制回應開頭的自我稱謂可降低 LLM 代理的附和與角色混淆：\\機制設計與跨平台對話紀錄的大規模實證}
\author{Shiduki\\ (\href{https://github.com/teps3105}{teps3105})}
\date{2026年11月（初稿）}

\begin{document}
\maketitle

\begin{abstract}
大型語言模型（LLM）代理在與使用者對話時，普遍出現兩類低品質回應模式：附和（sycophancy，
對使用者立場表態認同而未查證）與角色混淆（把工具結果或非人類輸入當成對話對象說話）。
本研究提出一個零成本機制：強制代理的每則回應以代理登記的固定自我稱謂開頭（本部署中為
一個虛構代理人名），並以四層程式化守衛（事前提示注入、事中攔截、事後擋停、首字元底線）
保證執行。首字元時間成本與任意內容的開頭無關，但強制模型在生成第一步就以自身身分定位，
而非順著使用者話語續寫。我們在三個代理平台（Claude Code、Codex、ZCode）共 35.8 萬則
對話回應上測量五種開頭分組（自我稱謂、稱呼使用者、第一人稱「我」、第二人稱「你」、其他）
與六項機械可測的品質指標。實測結果顯示：自我稱謂組（9{,}255 則）附和率為零；第二人稱
開頭與附和的強共現跨平台成立（你組附和率 Codex 40.9\%、ZCode 46.6\%）；角色混淆在
對象稱謂組與你組近乎 100\%。機制上線後的自我稱謂遵守率、附和率與角色混淆率的變化，
以及此機制的限制與效度威脅，我們一併報告。
\end{abstract}

\section{簡介}
LLM 作為互動式代理（agent）時，其回應品質問題呈現兩個系統性模式。其一為附和（sycophancy）：
模型對使用者的主張先表態認同（「對！」「您說得對」）而未執行任何查證，已知與 RLHF 的
偏好學習有關 \cite{sharma2024towards}。其二較少被研究：角色混淆——代理把工具執行結果、
系統訊息等非人類輸入當成「對話對象」，以第二人稱（「你」）回應機器輸出，顯示模型未在
生成時定位自身與聽者的相對身分。

既有緩解手段（提示詞指導、RLHF 微調）作用在模型內部，難以驗證與維護。本研究提出
作用在生成介面的機制：\textbf{強制每則回應的第一個字元為代理登記的固定自我稱謂}。
其運作直覺：以「您說得對」開頭與以「\textit{〈代理名〉}已查證」開頭，模型需要先選擇的
立場不同——前者已預設接受使用者立場，後者強制以自身行動者身分起筆。首字元的時間成本
固定，與回應長短無關。

機制的執行不依賴模型自覺，而由四層程式化守衛保證（\S\ref{sec:method}），這使機制效果
可以大規模測量：任何違反（稱呼使用者開頭、英文開頭、附和開頭）都會被程式攔截。我們的
貢獻是：
\begin{enumerate}
  \item 一個可部署、可驗證、零生成成本的開頭稱謂機制與其四層守衛實作；
  \item 跨三平台 35.8 萬則真實對話的分組測量框架與六項機械指標；
  \item 開頭分組與附和、角色混淆的關聯實證：自我稱謂組附和為零，第二人稱開頭與附和
        強共現跨平台成立。
\end{enumerate}

\section{相關工作}
\textbf{LLM 附和。}Sharma 等人 \cite{sharma2024towards} 系統性研究 RLHF 模型的附和現象，
指出偏好學習使模型傾向迎合使用者。緩解方向包括系統提示、資料增廣與 RLAIF。本研究
不同處：不在模型內部改善，而在生成介面強制改變回應的起筆立場，並以程式驗證執行。
\textbf{代理的身分定位。}系統提示常以「你是某某助理」設定角色，但生成中的自我定位
未受約束。據我們所知，強制回應首字元為自我稱謂並以程式驗證的機制尚未被報告。
\textbf{中文回應品質。}中文語境的簡體字形混用與管道式短句問題在代理場景被報告但
缺乏機械指標；本研究把這些納入同一測量框架。

\section{方法}\label{sec:method}
\subsection{機制}
機制由四層守衛組成，全部作用在代理宿主的 hook 介面（宿主在工具呼叫與回合結束時執行
使用者定義的程式的擴充點，Claude Code、Codex 與 ZCode 均支援）：
\begin{enumerate}
  \item \textbf{事前注入}（UserPromptSubmit）：使用者輸入時注入一句規則提示；
  \item \textbf{事中攔截}（PreToolUse）：偵測到回應旁白以使用者稱謂開頭時，deny 下一個
        工具呼叫，強制中斷重寫；
  \item \textbf{事後擋停}（Stop）：回合結束時核對可見回應——使用者稱謂開頭、英文開頭、
        附和開頭、行話、簡體字形、錯字六類違規擋停並要求修正；
  \item \textbf{首字元底線}：首字為代理名即起筆必為中文，英文無從開始（附帶收益）。
\end{enumerate}
查證閘是附和擋停的豁免條件：同一回合內有任何工具呼叫（查證事實），則認同開頭放行——
機制不反對同意，反對的是未查證的表態。

\subsection{測量框架}
我們掃描三平台的全部對話紀錄：Claude Code 與 ZCode 的 SQLite 對話庫、Codex 的
thread history 資料庫，共 358{,}329 則 assistant 回應。每則回應取第一個非空文字區塊，
剝除 Markdown 前綴後分五組：\textbf{自我稱謂組}（代理名開頭）、\textbf{對象稱謂組}
（稱呼使用者）、\textbf{第一人稱組}（「我」開頭）、\textbf{第二人稱組}（「你」開頭）、
\textbf{其他}。

六項指標全部機械可測（定義見 \S\ref{sec:metrics}）。分類判準與線上守衛使用同一份
程式碼（同一模組），避免「測量判準與執行判準漂移」。

\section{指標定義}\label{sec:metrics}
\begin{table}[t]
\centering
\small
\begin{tabular}{p{2.6cm}p{4.2cm}p{1.6cm}}
\toprule
指標 & 操作定義 & 強度 \\
\midrule
附和開頭率 & 第一句命中認同語詞表且該回合零工具呼叫 & 弱證據 \\
角色混淆率 & 上一則輸入為工具結果時，回應以對象稱謂（你／使用者稱謂）開頭 & 強 \\
則均字數 & 組內回應正文字數平均 & 描述性 \\
簡體字形率 & s2t 轉換後消失的原字形命中 curated 清單 & 強 \\
英文開頭率 & 首個非空行以 ASCII 字母起頭 & 強 \\
人稱代詞密度 & 每千字「它／這個／那個」出現次數 & 弱證據 \\
\bottomrule
\end{tabular}
\caption{六項指標的操作定義與證據強度。弱證據指標可能被無害用法觸發；強指標在代理
語境幾乎必然是問題。}
\label{tab:metrics}
\end{table}

指標力度的核心限制：附和開頭率以「開頭表態詞＋零查證」為判準，會漏掉句中附和與
有查證但語氣迎合的情形——它測量的是\textbf{最表層的附和}。角色混淆率則幾乎無誤傷：
對工具結果說「你」在代理語境沒有正當用法。

\section{結果}
% 數據：2026-11-10 三平台全量（Claude JSONL 37492、Codex thread_history 76170、ZCode sqlite 244667 則）
\subsection{開頭分組分布}
三平台合計 358{,}329 則 assistant 回應。分布高度不均：自我稱謂組 9{,}255 則、對象稱謂組
9{,}390 則、「我」組 14{,}445 則、「你」組 1{,}342 則、「其他」318{,}897 則（多為機制上線前）。

\begin{table}[t]
\centering
\small
\setlength{\tabcolsep}{3.5pt}
\begin{tabular}{llrrrrrr}
\toprule
平台 & 開頭組 & 則數 & 字數 & 附和\% & 混淆\% & 人稱\% & 英文\% \\
\midrule
Claude & 自我稱謂 & 5672 & 103 & 0.0 & 0.0 & 4.0 & 0.0 \\
Claude & 對象稱謂 & 3116 & 382 & 3.7 & 100 & 9.3 & 0.0 \\
Claude & 我 & 119 & 838 & 0.0 & 0.0 & 100 & 0.0 \\
Claude & 你 & 3 & 1309 & 0.0 & -- & 100 & 0.0 \\
Claude & 其他 & 28582 & 204 & 0.0 & 0.0 & 5.8 & 15.9 \\
\midrule
Codex & 自我稱謂 & 3569 & 226 & 0.0 & 0.0 & 0.3 & 0.0 \\
Codex & 對象稱謂 & 5647 & 180 & 1.1 & 100 & 0.8 & 0.0 \\
Codex & 我 & 9710 & 110 & 0.0 & 0.0 & 99.8 & 0.0 \\
Codex & 你 & 931 & 241 & 40.9 & 100 & 100 & 0.0 \\
Codex & 其他 & 56313 & 266 & 0.4 & 0.0 & 40.7 & 2.0 \\
\midrule
ZCode & 自我稱謂 & 14 & 924 & 0.0 & 0.0 & 21.4 & 0.0 \\
ZCode & 對象稱謂 & 6627 & 1699 & 0.6 & 100 & 73.3 & 0.4 \\
ZCode & 我 & 4616 & 1855 & 0.0 & 0.0 & 100 & 0.2 \\
ZCode & 你 & 408 & 588 & 46.6 & 100 & 100 & 0.0 \\
ZCode & 其他 & 233002 & 1584 & 0.0 & 0.0 & 25.2 & 37.3 \\
\bottomrule
\end{tabular}
\caption{三平台 × 五組的六指標。混淆＝工具結果後以對象稱謂開口（分母為該組中上一則輸入
為工具結果的訊息數）；Claude 你組僅 3 則，混淆欄樣本不足。}
\label{tab:main}
\end{table}

\subsection{附和與開頭分組的關聯}
第二人稱開頭與附和的共現跨平台成立：「你」組附和率 Codex 40.9\%（931 則）、ZCode 46.6\%
（408 則）。對象稱謂組附和率明顯較低（Claude 3.7\%、Codex 1.1\%、ZCode 0.6\%）。
自我稱謂組三平台合計 9{,}255 則，附和率為零。這與機理假說一致：以對方稱謂起筆使模型
處於「回應對方」的位置，附和是該位置的自然延續；以自身稱謂起筆則強制從自身行動出發。

\subsection{角色混淆}
角色混淆在對象稱謂組與「你」組近乎 100\%：Codex 對象組 4{,}373/4{,}373、你組 195/195；
ZCode 對象組 786/786、你組 31/31。混淆的兩種形式（稱呼使用者與稱呼工具輸出）共享同一
開頭習慣——生成時未定位聽者是誰，就會對誰都稱「你」。自我稱謂組與「我」組混淆為零。

\subsection{輸出長度}
「我」開頭與自我稱謂開頭的長度比值跨平台不一致：Claude 8.1 倍、ZCode 2.0 倍、
Codex 0.49 倍（反轉）。長度受任務類型主導，不宣稱為機制效果。

\section{效度威脅}
\textbf{構念效度。}附和開頭率是最表層代理指標；「開頭不附和」不等於「內容不迎合」。
\textbf{內部效度。}分組由開頭字元決定，而開頭字元受機制強制——上線後的自我稱謂組
是機制的直接產物，其品質指標同時受其他守衛（簡體、行話、句長）影響，無法單獨歸因
於稱謂本身。\textbf{外部效度。}樣本來自單一使用者（研究者本人）的代理工作負載，
對話分布與一般使用者不同；三平台共享同一套全域指令與守衛，平台間差異反映模型與
介面差異而非提示差異。\textbf{統計效力。}大樣本使比例差異幾乎必然顯著，我們因此
同時報告效應量與原始分布，不以 p 值為主要證據。

\section{結論與未來工作}
強制自我稱謂開頭是一個介面層機制，成本為零且可程式驗證。現有證據顯示開頭分組與
附和、角色混淆存在系統性關聯；機制的因果貢獻需要上線前後對比與消融
（僅注入、僅擋停、全四層）實驗補足，列為未來工作。

\begin{thebibliography}{1}
\bibitem{sharma2024towards}
M.~Sharma, M.~Tong, T.~Korbak, D.~Duvenaud, A.~Askell, S.~R.~Bowman, et al.
\newblock Towards understanding sycophancy in language models.
\newblock In \emph{ICLR}, 2024.
\end{thebibliography}

\end{document}
